% ============================================================
% Core packages
% ============================================================
\usepackage{subfiles}        % Modular compilation
\usepackage{pdfpages}        % Include external PDFs
\usepackage{lipsum}          % Dummy text (remove for final version)
\usepackage{microtype}
\usepackage{etoolbox}

% Language and hyphenation
\usepackage[main=english, spanish, danish, greek, french]{babel}
\frenchsetup{AutoSpacePunctuation=false} % no automatic space before : ; ! ?
\usepackage{csquotes}

% Figures and Tables
\usepackage{float}
\usepackage{booktabs}
\usepackage[export]{adjustbox}
\usepackage{graphicx}
\usepackage[figurename=Fig.]{caption}
\captionsetup{font=small}

% Other scientific packages
\usepackage{chemmacros}
\chemsetup{greek=upgreek}
\usepackage{siunitx}
\usepackage{nicematrix}
\usepackage{physics}

% TikZ
\usepackage{tikz}

% Misc
\usepackage{ragged2e}
\usepackage{indentfirst}
\usepackage{footnotehyper}
\usepackage{multicol}

\usepackage[acronym, toc, nomain]{glossaries}
\makeglossaries
% Disable glossaries' own per-entry hyperlinks (which mis-anchor in memoir's
% boxed item labels).  Instead, \glstextformat wraps every in-text acronym in
% a plain \hyperref link that goes to the \label{abbreviations} on the list
% page — reliable, no colour, still clickable.
\glsdisablehyper
\makeatletter
\newcommand{\thesisacronymlink}[1]{%
	\@ifundefined{r@abbreviations}{#1}{\hyperref[abbreviations]{\textcolor{black}{#1}}}%
}
\makeatother
\renewcommand{\glstextformat}[1]{\thesisacronymlink{#1}}

% Two-column, tight-spacing glossary style for the List of Abbreviations:
% clones the stock "list" style but wraps the entries (not the chapter
% heading printed by \printglossary) in a multicols{2} block, and collapses
% the description-list spacing since most entries are shorter than a line.
\newglossarystyle{acrtwocol}{%
	\setglossarystyle{list}%
	\renewenvironment{theglossary}%
		{\glslistinit
		 \begin{multicols}{2}%
		 \setlength{\parskip}{0pt}%
		 \setlength{\itemsep}{0pt}%
		 \setlength{\parsep}{0pt}%
		 \setlength{\topsep}{0pt}%
		 \setlength{\partopsep}{0pt}%
		 \begin{description}}%
		{\end{description}\end{multicols}}%
}
